\paragraph{What the evidence says.}
Against the seven hypotheses we pre-specified (Table~\ref{tab:hyp}, Appendix), H1, H3 and H5 are supported, H2, H4 and H6 partly, and H7 not as a general law.
Regression adjustment for the GATE is a bias--variance bargain whose terms are set by the design and the graph, not by the sample size alone. When the exposure model is right it is the most accurate of the estimators we compared once $n$ and the interference are large enough (in the default SBM setting, $n\gtrsim2{,}000$ and interference at least as strong as the direct effect; the break-even moves with degree and confounding); under Bernoulli$(0.5)$ its median relative RMSE over the six main-grid graphs is then \MainOneOls{} (DM: \MainDM). It is not a free lunch: its variance grows with degree and with community-level baseline confounding, and a difference in means is better with weak interference, small $n$, or after cluster randomisation. When the exposure model is wrong, the error is a bias that more data do not remove; we can compute its first-order size in the linear family (P1) and observe it for nonlinear, two-hop, degree-heterogeneous and miswired responses.
\paragraph{Practical reading.} (i) Fit the simplest exposure model, but check the design's exposure support: bias is largest where $f$ is far from the design's mean exposure. (ii) Prefer balanced designs; at $p=0.1$ nonlinear bias and variance both grow. (iii) Do not rely on high-capacity learners to escape misspecification. (iv) Treat nominal $95\%$ intervals, cluster-robust ones included, as optimistic. These are readings of simulations.
\paragraph{Limitations.}
\emph{Synthetic outcomes.} The graphs are real or generated, but every outcome is simulated from mechanisms we chose; we did not test the estimators on any experiment with real outcomes, and our mechanisms (linear, saturating, power, threshold, two-hop, degree-heterogeneous) are a small family. The relative scale (e.g.\ DM at $0.5$) is a function of the ratio $\rho$ we chose. \emph{Replication.} Synthetic cells use 5--8 graph draws; each real graph is one graph, and its baseline draw, which the bootstrap does not cover, matters (Sec.~\ref{sec:struct}). Sweeps vary one factor at a time. The model-class comparison (Fig.~\ref{fig:model}) uses 4 graph draws $\times$ 80 assignments, and a ridge/boosting evaluation on the full main grid was not completed (an attempt was stopped; Appendix~\ref{app:setup}). \emph{Analytics.} P1--P3 are first-order large-$n$ arguments for the linear-exposure family, checked numerically, not proved, and do not cover boosted or cubic fits. \emph{Inference.} Louvain communities are both the cluster design and, in some cells, the variance clusters, which favours the cluster-robust interval; we studied no corrections. \emph{Not done.} Multi-$p$ and staggered designs (H4), the dependence of coverage on $n$ under misspecification (H6), unknown-interference-radius selection, and optimal-design recommendations. The H7 diagnostic did not generalise, so we offer no data-driven warning signal.
